\chapter{Compact Groups}

\begin{center}\itshape
Averaging over a compact group with its Haar measure makes representations unitary and characters orthonormal, as for finite groups. The Peter--Weyl theorem says that matrix coefficients are complete, like the exponentials of Fourier series. For the
unitary group, Weyl's integration formula then finds every irreducible character: they are the Schur polynomials. And the characters compute moments of random unitary matrices, which turn out to count permutations by their longest increasing subsequence.
\end{center}

In 1924 Issai Schur took Hurwitz's invariant integral over the orthogonal group (Tome~8, Chapter~8) and showed that with it, Frobenius' theory of characters works for rotation groups as it does for finite groups: sums over the group become integrals.
Hermann Weyl, in Zurich, read the paper and saw much further. In three long articles of 1925 and 1926 he found the characters of all the compact semisimple Lie groups, with his integration formula and a character formula that bears his name; he also
used integration over a compact group to prove that representations of complex Lie algebras such as $\mathfrak{sl}_n(\C)$ are sums of irreducibles, a method now called the ``unitarian trick''. In 1927, with his student Fritz Peter, he proved
that the matrix coefficients of the irreducible representations of a compact group are complete, as the exponentials $e^{ik\theta}$ are complete on the circle. Peter later became the headmaster of a secondary school; his name stays attached to one of the
pillars of harmonic analysis.

In 1939 Weyl described algebra and topology as fighting over every field of mathematics: ``In these days the angel of topology and the devil of abstract algebra fight for the soul of each individual mathematical domain.'' Compact groups are a field where they
make peace. This chapter takes the analytic side (Haar measure, Stone--Weierstrass, $L^2$) and the algebraic side (characters, alternating polynomials, Young diagrams) and shows them computing the same numbers. It ends where Tome~8 ended, with random unitary
matrices, and with a surprise linking them to the longest increasing subsequences of Chapter~9.

Throughout, $G$ is a compact matrix Lie group (Tome~8, Chapter~7), a closed subgroup of some $\mathrm{GL}_N(\C)$ which is compact, such as $U(n)$, $SU(n)$ or $SO(n)$, and $\mu$ is its Haar probability measure, invariant on both sides (Tome~8, Theorem~8.3 and
Corollary~8.2). A \emph{representation} of $G$ is a continuous homomorphism $\rho:G\to\mathrm{GL}(V)$ with $V$ a finite-dimensional complex vector space.

\section{Representations of compact groups}

Everything in Tome~2, Chapter~11, that used the average $\frac1{|G|}\sum_g$ over a finite group carries over with $\int_G\dd\mu(g)$ in its place: the invariance $\int f(hg)\dd\mu(g)=\int f(g)\dd\mu(g)$ is all that was used.

\begin{proposition}\label{prop:unitarytrick}
Every representation $\rho$ of $G$ is unitary for some inner product on $V$, and is a direct sum of irreducible representations.
\end{proposition}

\begin{proof}
Start from any inner product $\langle\cdot,\cdot\rangle$ and set $\langle v,w\rangle_G=\int_G\langle\rho(g)v,\rho(g)w\rangle\dd\mu(g)$, the integral of a continuous function. It is an inner product (positive, since the integrand is positive for $v\ne0$ and $\mu$ is a
probability), and $\langle\rho(h)v,\rho(h)w\rangle_G=\langle v,w\rangle_G$ by invariance of $\mu$ under $g\mapsto gh$. For this inner product, the orthogonal complement of an invariant subspace is invariant, and induction on the dimension splits $V$ into irreducibles.
\end{proof}

From now on representations are unitary, in orthonormal bases, so that $\rho(g^{-1})=\rho(g)^*$, and their matrix entries $\rho_{ij}(g)$ are continuous functions on $G$, the \emph{matrix coefficients}. Schur's lemma holds as for finite groups (its proof uses only
linear algebra), and so does the following.

\begin{theorem}[Schur orthogonality]\label{thm:schurorth}
Let $\rho,\sigma$ be irreducible unitary representations, of dimensions $d_\rho,d_\sigma$. Then
\[
\int_G\rho_{ij}(g)\,\overline{\sigma_{kl}(g)}\dd\mu(g)=\begin{cases}\frac1{d_\rho}\delta_{ik}\delta_{jl}&\text{if }\rho=\sigma,\\0&\text{if }\rho\not\cong\sigma.\end{cases}
\]
Consequently the characters $\chi_\rho=\operatorname{tr}\rho$ of irreducible representations are orthonormal in $L^2(G,\mu)$, and two representations with the same character are isomorphic.
\end{theorem}

\begin{proof}
For a linear map $A:V_\sigma\to V_\rho$, let $\tilde A=\int_G\rho(g)A\sigma(g)^{-1}\dd\mu(g)$ (entrywise integrals). Substituting $g\mapsto hg$ shows $\rho(h)\tilde A\sigma(h)^{-1}=\tilde A$: $\tilde A$ is a module map. By Schur's lemma $\tilde A=0$ if $\rho\not\cong\sigma$,
and if $\rho=\sigma$ it is a scalar, which is $\frac{\operatorname{tr}A}{d_\rho}$ since $\operatorname{tr}\tilde A=\operatorname{tr}A$. Take for $A$ the matrix unit $E_{jl}$. Then $\tilde A_{ik}=\int\rho_{ij}(g)\,\sigma(g^{-1})_{lk}\dd\mu$, and $\sigma(g^{-1})_{lk}=\overline{\sigma_{kl}(g)}$
by unitarity; while $\operatorname{tr}E_{jl}=\delta_{jl}$. This is the formula. Summing it over $i=j$ and $k=l$ gives $\langle\chi_\rho,\chi_\sigma\rangle=\delta_{\rho\sigma}$. Finally, by Proposition~\ref{prop:unitarytrick} a representation is $\bigoplus m_\rho\rho$, its
character is $\sum m_\rho\chi_\rho$, and $m_\rho=\langle\chi,\chi_\rho\rangle$ is read off the character.
\end{proof}

For the circle $U(1)$, the irreducible representations are the characters $z\mapsto z^k$, $k\in\Z$ (Tome~7, Theorem~8.11), and Schur orthogonality is $\frac1{2\pi}\int_0^{2\pi}e^{ik\theta}e^{-il\theta}\dd\theta=\delta_{kl}$. The question that Fourier answered on the
circle, whether these functions are \emph{complete}, is the subject of the next section.

\section{The Peter--Weyl theorem}

Let $\widehat G$ be a set of representatives of the irreducible representations of $G$, and $\mathcal R(G)$ the span of all matrix coefficients of all representations: the \emph{representative functions}.

\begin{theorem}[Peter--Weyl, 1927]\label{thm:peterweyl}
\begin{enumerate}[label=(\alph*)]
\item $\mathcal R(G)$ is dense in $C(G)$ for the uniform norm.
\item The functions $\sqrt{d_\rho}\,\rho_{ij}$, for $\rho\in\widehat G$ and $1\le i,j\le d_\rho$, form an orthonormal basis of $L^2(G,\mu)$. So every $f\in L^2$ is the $L^2$-convergent sum of its \emph{Fourier series}
\[
f=\sum_{\rho\in\widehat G}d_\rho\operatorname{tr}\bigl(\hat f(\rho)\rho\bigr),\qquad\hat f(\rho)=\int_Gf(g)\rho(g)^*\dd\mu(g),
\]
with the Plancherel identity
\[
\int_G|f|^2\dd\mu=\sum_\rho d_\rho\operatorname{tr}\bigl(\hat f(\rho)\hat f(\rho)^*\bigr).
\]
\item The characters $\chi_\rho$, $\rho\in\widehat G$, form an orthonormal basis of the closed subspace of $L^2(G,\mu)$ of class functions, those with $f(hgh^{-1})=f(g)$ for all $h$ and almost every $g$.
\end{enumerate}
\end{theorem}

\begin{proof}
(a) $\mathcal R(G)$ is a subspace by definition, and a subalgebra of $C(G)$: the product $\rho_{ij}\sigma_{kl}$ is the coefficient of $\rho\otimes\sigma$ in the position $((i,k),(j,l))$. It contains the constants (the trivial representation), and is stable under conjugation: $\overline{\rho_{ij}}$ is a coefficient of the conjugate representation
$g\mapsto\overline{\rho(g)}$, which is continuous and multiplicative. It separates points: $G\subseteq\mathrm{GL}_N(\C)$, and the entries $g\mapsto g_{ab}$ are the coefficients of the defining representation, which distinguish any two distinct matrices. By the
Stone--Weierstrass theorem (Tome~6, Theorem~4.19(b)), $\mathcal R(G)$ is dense in $C(G)$.

(b) Orthonormality is Theorem~\ref{thm:schurorth}. The span of these functions is $\mathcal R(G)$: a representation is a sum of irreducibles (Proposition~\ref{prop:unitarytrick}), so in a suitable basis its coefficients are coefficients of irreducibles or $0$,
and a change of basis replaces coefficients by linear combinations of them. Now $C(G)$ is dense in $L^2(\mu)$: simple functions are dense (Tome~6, Chapter~9), and the indicator of a Borel set $E$ is an $L^2$ limit of continuous functions, because the finite
Borel measure $\mu$ on the compact metric space $G$ is regular (for $\varepsilon>0$ there are $F\subseteq E\subseteq O$, $F$ closed, $O$ open, with $\mu(O\setminus F)<\varepsilon$), and a continuous function $0\le\phi\le1$ that is $1$ on $F$ and $0$ off $O$ exists (Urysohn,
or $\phi(x)=\frac{d(x,O^c)}{d(x,O^c)+d(x,F)}$). Uniform approximation implies $L^2$ approximation for a probability measure, so $\mathcal R(G)$ is dense in $L^2$, and an orthonormal family with dense span is a basis. The coefficient of $f$ on $\sqrt{d_\rho}\rho_{ij}$ is
$\sqrt{d_\rho}\int f\overline{\rho_{ij}}\dd\mu=\sqrt{d_\rho}\,\hat f(\rho)_{ji}$, and the expansion and Parseval's identity in this basis are the stated formulas.

(c) For $\rho\in\widehat G$, let $E_\rho$ be the span of the $\rho_{ij}$, and $P_\rho$ the orthogonal projection onto it. For $h\in G$, the map $(c_hf)(g)=f(h^{-1}gh)$ is unitary on $L^2$ (invariance of $\mu$ on both sides) and maps $E_\rho$ to itself, since
$\rho(h^{-1}gh)=\rho(h)^{-1}\rho(g)\rho(h)$. So $P_\rho$ commutes with $c_h$, and if $f$ is a class function, $P_\rho f$ is a class function in $E_\rho$. Write $P_\rho f(g)=\operatorname{tr}(A\rho(g))$ for a matrix $A$ (every element of $E_\rho$ has this form,
and $A$ is unique since the $\rho_{ij}$ are independent). Invariance under $c_h$ reads $\operatorname{tr}(\rho(h)A\rho(h)^{-1}\rho(g))=\operatorname{tr}(A\rho(g))$, so $\rho(h)A\rho(h)^{-1}=A$ for all $h$, and $A$ is a scalar $a$ by Schur's lemma: $P_\rho f=a\chi_\rho$, with
$a=\langle f,\chi_\rho\rangle$. By (b), $f=\sum_\rho P_\rho f=\sum_\rho\langle f,\chi_\rho\rangle\chi_\rho$, and the characters, orthonormal by Theorem~\ref{thm:schurorth}, are a basis of the class functions.
\end{proof}

For $U(1)$, part (b) is the completeness of Fourier series, here proved with Stone--Weierstrass instead of Fejér's kernel; part (a) for $SO(3)$ says that spherical harmonics approximate every continuous function on the sphere (Exercise~\ref{ex:su2}
shows the corresponding fact for $SU(2)$). Peter and Weyl proved the theorem for compact Lie groups, and after Haar's theorem of 1933 the same proof works for every compact group; it uses the spectral theorem for compact self-adjoint operators (Tome~7, Theorem~4.4) in place of Stone--Weierstrass: convolution with a continuous function
is a compact operator commuting with translations, and its finite-dimensional eigenspaces are representations. One consequence is that every compact Lie group is a matrix group, so that our restriction loses nothing.

\begin{think}
Where did compactness enter? Without a finite invariant measure, averaging an inner product is impossible, and indeed the representation $t\mapsto\begin{pmatrix}1&t\\0&1\end{pmatrix}$ of $(\R,+)$ is not a sum of irreducibles: the line $\C e_1$ is
invariant and has no invariant complement. Is there a non-compact group all of whose representations are sums of irreducibles?
\end{think}

\section{The irreducible representations of \texorpdfstring{$U(n)$}{U(n)}}

Let $T\subseteq U(n)$ be the diagonal torus, $t=\mathrm{diag}(t_1,\dots,t_n)$. A character of $U(n)$ is a class function, determined by its values on $T$ (every unitary matrix is conjugate to an element of $T$, Tome~8, Theorem~7.9), where it is
symmetric in $t_1,\dots,t_n$ (permutation matrices conjugate $T$ to itself, permuting the entries). Weyl's integration formula (Tome~8, Theorem~8.8) computes inner products of class functions on $T$ alone. So the problem is to find the right symmetric functions.

For $\mu=(\mu_1,\dots,\mu_n)\in\Z^n$ write $t^\mu=t_1^{\mu_1}\cdots t_n^{\mu_n}$, and let
\[
a_\mu(t)=\det\bigl(t_k^{\mu_j}\bigr)_{1\le j,k\le n}=\sum_{w\in\mathfrak S_n}\operatorname{sgn}(w)\,t_{w(1)}^{\mu_1}\cdots t_{w(n)}^{\mu_n},
\]
an \emph{alternant}: it changes sign when two variables are exchanged. Let $\delta=(n-1,n-2,\dots,1,0)$; then $a_\delta=\prod_{j<k}(t_j-t_k)$, the Vandermonde determinant (Tome~2). A \emph{dominant weight} is $\lambda\in\Z^n$ with $\lambda_1\ge\lambda_2\ge\dots\ge\lambda_n$;
then $\lambda+\delta$ is strictly decreasing.

\begin{definition}\label{def:schurpoly}
The \emph{Schur function} of a dominant weight $\lambda$ is $s_\lambda=a_{\lambda+\delta}/a_\delta$.
\end{definition}

\begin{lemma}\label{lem:alternants}
\begin{enumerate}[label=(\alph*)]
\item $s_\lambda$ is a symmetric Laurent polynomial with integer coefficients; its largest monomial in the lexicographic order is $t^\lambda$, with coefficient $1$.
\item Every antisymmetric Laurent polynomial is a unique finite combination $\sum_\mu c_\mu a_\mu$ over strictly decreasing $\mu$, with $c_\mu$ its coefficient of $t^\mu$.
\item For strictly decreasing $\mu,\nu$,
\[
\frac1{n!}\int_T a_\mu\overline{a_\nu}\,\frac{\dd\theta}{(2\pi)^n}=\delta_{\mu\nu},
\]
where $t_j=e^{i\theta_j}$ and $\dd\theta=\dd\theta_1\cdots\dd\theta_n$.
\end{enumerate}
\end{lemma}

\begin{proof}
(b) Let $F=\sum_\mu c_\mu t^\mu$ be antisymmetric. Exchanging $t_j$ and $t_k$ shows that $c_\mu=0$ when $\mu_j=\mu_k$, and in general $c_{w\mu}=\operatorname{sgn}(w)c_\mu$, where $w$ permutes the entries. Grouping the monomials by the decreasing
rearrangement of their exponents gives $F=\sum c_\mu a_\mu$; uniqueness holds because distinct strictly decreasing $\mu$ give alternants with disjoint sets of monomials.

(a) Multiplying by a power of $t_1\cdots t_n$, we may assume that $a_{\lambda+\delta}$ is a polynomial. It vanishes when $t_j=t_k$, so, as a polynomial in $t_j$ with coefficients in the other variables, it is divisible by $t_j-t_k$. The factors $t_j-t_k$ are pairwise
non-associate irreducibles in the unique factorisation domain $\Z[t_1,\dots,t_n]$, so $a_\delta$ divides $a_{\lambda+\delta}$, with a quotient that has integer coefficients and is symmetric (a ratio of two antisymmetric functions). Order $\Z^n$ lexicographically; this
order is compatible with addition, so the largest monomial of a product is the product of the largest monomials. The largest monomial of $a_\nu$, for $\nu$ strictly decreasing, is $t^\nu$ (from $w=1$), with coefficient $1$. So the largest monomial of $s_\lambda$
is $t^{\lambda+\delta}/t^\delta=t^\lambda$, with coefficient $1$.

(c) The integral of $t^\alpha\overline{t^\beta}$ over $T$ is $1$ if $\alpha=\beta$ and $0$ otherwise. The alternant $a_\mu$ is a sum of $n!$ distinct monomials with coefficients $\pm1$, disjoint from those of $a_\nu$ when $\nu\ne\mu$.
\end{proof}

\begin{theorem}[Weyl's character formula for $U(n)$, 1925]\label{thm:weylchar}
For every dominant weight $\lambda$ there is an irreducible representation $V_\lambda$ of $U(n)$ whose character on $T$ is $s_\lambda$, and every irreducible representation is isomorphic to exactly one $V_\lambda$. Its dimension is
\[
\dim V_\lambda=s_\lambda(1,\dots,1)=\prod_{1\le j<k\le n}\frac{\lambda_j-\lambda_k+k-j}{k-j}.
\]
\end{theorem}

\begin{proof}
\emph{Weights.} Let $\rho$ be a representation. The operators $\rho(t)$, $t\in T$, are unitary and commute, so they have a common orthonormal basis of eigenvectors. On each eigenline $T$ acts by a continuous homomorphism $T\to\C^\times$, with image a compact subgroup, hence in
the unit circle; restricted to each factor $U(1)$ it is $z\mapsto z^{k}$ (Tome~7, Theorem~8.11), so it is $t\mapsto t^\mu$ for some $\mu\in\Z^n$, a \emph{weight} of $\rho$. Hence $\chi_\rho(t)=\sum_\mu m_\mu t^\mu$, a symmetric Laurent polynomial whose
coefficients $m_\mu\ge0$ are the multiplicities of the weights.

\emph{Every irreducible character is a Schur function.} Let $\rho$ be irreducible. Then $\chi_\rho a_\delta$ is antisymmetric with integer coefficients, so $\chi_\rho a_\delta=\sum_\nu c_\nu a_\nu$ with $c_\nu\in\Z$, by Lemma~\ref{lem:alternants}(b). By Weyl's integration formula,
in which $\prod_{j<k}|t_j-t_k|^2=|a_\delta|^2$, and by Lemma~\ref{lem:alternants}(c),
\[
1=\int_{U(n)}|\chi_\rho|^2\dd\mu=\frac1{n!}\int_T|\chi_\rho a_\delta|^2\frac{\dd\theta}{(2\pi)^n}=\sum_\nu c_\nu^2.
\]
So exactly one $c_\nu$ is non-zero, and it is $\pm1$: $\chi_\rho=\pm a_\nu/a_\delta=\pm s_\lambda$ with $\lambda=\nu-\delta$ dominant. The coefficient of $t^\lambda$ in $s_\lambda$ is $1$ (Lemma~\ref{lem:alternants}(a)), and in $\chi_\rho$ it is a multiplicity $m_\lambda\ge0$; so the sign is
$+$. Non-isomorphic irreducibles have orthogonal, hence different, characters, so they correspond to different $\lambda$.

\emph{Every Schur function is a character.} Let $\lambda$ be dominant, and suppose no irreducible representation has character $s_\lambda$. A symmetric Laurent polynomial is a polynomial in the elementary symmetric functions $e_1,\dots,e_n$ and $e_n^{-1}$
(Tome~0, Theorem~3.11), and these are, up to sign, the coefficients of the characteristic polynomial; so $f(U)=s_\lambda(\text{eigenvalues of }U)$ is a continuous class function on $U(n)$. By Weyl's formula and
Lemma~\ref{lem:alternants}(c), $\langle f,\chi_\rho\rangle=\frac1{n!}\int_Ta_{\lambda+\delta}\overline{a_{\lambda'+\delta}}\frac{\dd\theta}{(2\pi)^n}=0$ for every irreducible $\rho$, whose character is some $s_{\lambda'}$ with $\lambda'\ne\lambda$, while $\langle f,f\rangle=1$. This contradicts
Theorem~\ref{thm:peterweyl}(c).

\emph{Dimension.} The dimension is $\chi(I)=s_\lambda(1,\dots,1)$, a quotient $0/0$, which we compute as a limit along $t=(q^{n-1},q^{n-2},\dots,1)$, $q\to1$. With $\nu=\lambda+\delta$, $a_\nu(q^{n-1},\dots,q^0)=\det\bigl(q^{(n-k)\nu_j}\bigr)$ is a Vandermonde
determinant in the variables $x_j=q^{\nu_j}$ (with exponents $n-k$), so it equals $\prod_{j<k}(q^{\nu_j}-q^{\nu_k})$; similarly for $\delta$. Hence
\[
s_\lambda(q^{n-1},\dots,1)=\prod_{j<k}\frac{q^{\nu_j}-q^{\nu_k}}{q^{\delta_j}-q^{\delta_k}}\xrightarrow[q\to1]{}\prod_{j<k}\frac{\nu_j-\nu_k}{\delta_j-\delta_k}=\prod_{j<k}\frac{\lambda_j-\lambda_k+k-j}{k-j},
\]
since $\frac{q^a-q^b}{q-1}\to a-b$.
\end{proof}

\begin{figure}[ht]
\centering
\fig{ch11-weights}
\caption{Weight diagrams of two representations of $U(3)$, drawn in the plane $\mu_1+\mu_2+\mu_3=\text{constant}$. Left: $V_{(2,1,0)}$, of dimension $8$: the six permutations of $(2,1,0)$, and the weight $(1,1,1)$ with multiplicity $2$ (the circled centre); restricted to $SU(3)$,
this is Gell-Mann's ``eightfold way'' of 1961, in which the eight lightest baryons fill the weights. Right: $V_{(3,0,0)}=\operatorname{Sym}^3\C^3$, of dimension $10$, whose tenth weight, at a corner of the triangle, was a particle not yet observed; Gell-Mann
predicted it in 1962, and the $\Omega^-$ was found at Brookhaven in 1964.}
\label{fig:weights}
\end{figure}

For $\lambda=(1,0,\dots,0)$, $s_\lambda=t_1+\dots+t_n$: the defining representation $\C^n$. For $\lambda=(1,\dots,1,0,\dots,0)$ with $k$ ones, $s_\lambda=e_k$, the elementary symmetric function: the exterior power $\Lambda^k\C^n$. For $\lambda=(k,0,\dots,0)$, $s_\lambda$ is
the sum of all monomials of degree $k$: $\operatorname{Sym}^k\C^n$. For $\lambda=(m,\dots,m)$, $s_\lambda=(t_1\cdots t_n)^m$: the one-dimensional $\det^m$. Twisting by $\det$ adds $1$ to every $\lambda_j$, so the essential case is $\lambda_n\ge0$, a partition with at most
$n$ rows. The dimension formula then has a version in terms of Young diagrams, $\dim V_\lambda=\prod_{(i,j)\in\lambda}\frac{n+j-i}{h_{ij}}$ (the \emph{hook-content formula}, of the same family as Theorem~\ref{thm:hook}).

The symmetric groups and the unitary groups are tied more closely still. $\mathfrak S_k$ permutes the factors of $(\C^n)^{\otimes k}$ and $U(n)$ acts on all factors at once, and these actions commute. Schur's thesis of 1901 showed that, as a representation of
$\mathfrak S_k\times U(n)$,
\[
(\C^n)^{\otimes k}\cong\bigoplus_{\lambda\vdash k,\ \ell(\lambda)\le n}S^\lambda\otimes V_\lambda,
\]
where $\ell(\lambda)$ is the number of rows: \emph{Schur--Weyl duality}. We will not prove it, but its trace on $U(n)$ is proved below (Lemma~\ref{lem:pieri} and Theorem~\ref{thm:traces}): $V_\lambda$ appears in $(\C^n)^{\otimes k}$ exactly $f^\lambda=\dim S^\lambda$ times.

\section{Traces of random unitary matrices}

Let $U$ be a Haar random element of $U(n)$. Its eigenvalues are spread around the circle and repel each other (Tome~8, Chapter~8), so they are spread more evenly than independent points. Here is a precise measurement.

\begin{proposition}\label{prop:powertraces}
For integers $j,k\ge1$, $\mathbb E[\operatorname{tr}U^j]=0$ and
\[
\mathbb E\bigl[\operatorname{tr}U^j\,\overline{\operatorname{tr}U^k}\bigr]=\delta_{jk}\min(k,n).
\]
\end{proposition}

\begin{proof}
If $z\in U(1)$, then $zU$ is again Haar distributed, so $\mathbb E[\operatorname{tr}U^j\overline{\operatorname{tr}U^k}]=z^{j-k}\,\mathbb E[\operatorname{tr}U^j\overline{\operatorname{tr}U^k}]$ for all $z$, which forces the expectation to vanish unless $j=k$; similarly $\mathbb E[\operatorname{tr}U^j]=0$.
For $j=k$, $\operatorname{tr}U^k=p_k(t)=\sum_it_i^k$ on $T$, the power sum. Since $p_k$ is symmetric,
\[
p_k\,a_\delta=\sum_{w}\operatorname{sgn}(w)\,w\bigl(t^\delta\bigr)\,w(p_k)=\sum_{i=1}^n\sum_w\operatorname{sgn}(w)\,w\bigl(t^{\delta+ke_i}\bigr)=\sum_{i=1}^na_{\delta+ke_i},
\]
where $w$ acts by permuting the variables and $e_i$ is the $i$-th unit vector. The vector $\delta+ke_i$ has the entries $n-1,\dots,0$, except that $n-i$ is raised to $n-i+k$. If $k\ge i$, this new entry is at least $n$, larger than all others, and
$a_{\delta+ke_i}=\pm a_\nu$ with $\nu$ the decreasing rearrangement; if $k<i$, it equals the entry $n-i+k$ of position $i-k$, and $a_{\delta+ke_i}=0$. The surviving $\nu$, one for each $i\le\min(k,n)$, are distinct (each is $\delta$ with the entry $n-i$ replaced). By
Weyl's formula and Lemma~\ref{lem:alternants}(c), $\mathbb E|\operatorname{tr}U^k|^2=\frac1{n!}\int_T|p_ka_\delta|^2=\min(k,n)$.
\end{proof}

Compare with $n$ independent uniform points $z_i$ on the circle: $\mathbb E|\sum z_i^k|^2=n$ for every $k$. The eigenvalues of a random unitary matrix are much more rigid: for small $k$, the sum $\sum t_i^k$ has variance $k$ instead of $n$. Diaconis and Shahshahani
proved in 1994 that, as $n\to\infty$, $\operatorname{tr}U,\operatorname{tr}U^2,\dots,\operatorname{tr}U^k$ become independent complex Gaussians of variances $1,2,\dots,k$. Here is the first coordinate, through a remarkable formula.

\begin{lemma}[Pieri's rule for one box]\label{lem:pieri}
For a dominant weight $\lambda$, $p_1s_\lambda=\sum_\nu s_\nu$, the sum over the dominant weights $\nu$ obtained from $\lambda$ by adding $1$ to one entry.
\end{lemma}

\begin{proof}
As in the proof of Proposition~\ref{prop:powertraces} with $k=1$, $p_1a_{\lambda+\delta}=\sum_ia_{\lambda+\delta+e_i}$. If $i>1$ and $\lambda_{i-1}=\lambda_i$, the entry $\lambda_i+n-i+1$ equals the entry in position $i-1$, and the term vanishes. Otherwise
$\lambda+e_i$ is dominant and the term is $a_{(\lambda+e_i)+\delta}$. Divide by $a_\delta$.
\end{proof}

For a partition $\lambda$ with at most $n$ rows (padded with zeros), adding $1$ to an entry while staying dominant is adding a box to the Young diagram, in a row $\le n$.

\begin{theorem}[Diaconis--Shahshahani 1994, Rains 1998]\label{thm:traces}
For every $a\ge0$, $(\operatorname{tr}U)^a=\sum_{\lambda\vdash a,\ \ell(\lambda)\le n}f^\lambda\chi_{V_\lambda}(U)$, and
\[
\mathbb E\bigl|\operatorname{tr}U\bigr|^{2a}=\#\bigl\{\sigma\in\mathfrak S_a:\ \text{the longest increasing subsequence of $\sigma$ has length}\le n\bigr\}.
\]
In particular $\mathbb E|\operatorname{tr}U|^{2a}=a!$ whenever $a\le n$.
\end{theorem}

\begin{proof}
On $T$, $\operatorname{tr}U=p_1$. Applying Lemma~\ref{lem:pieri} $a$ times starting from $s_{(0,\dots,0)}=1$, $p_1^a=\sum_\lambda N_\lambda s_\lambda$, where $N_\lambda$ is the number of ways of building the diagram $\lambda$ one box at a time, staying a Young diagram with
at most $n$ rows. Recording in each box the time at which it was added identifies these ways with the standard tableaux of shape $\lambda$, so $N_\lambda=f^\lambda$; this is the first formula, by Theorem~\ref{thm:weylchar}. The characters $s_\lambda$ are orthonormal,
so $\mathbb E|\operatorname{tr}U|^{2a}=\sum_{\lambda\vdash a,\,\ell(\lambda)\le n}(f^\lambda)^2$. Conjugation of diagrams swaps rows and columns and preserves $f^\lambda$ (Corollary~\ref{cor:conjugate}), so this is also the sum of $(f^\lambda)^2$ over $\lambda\vdash a$ with first
row $\lambda_1\le n$. By the Robinson--Schensted correspondence (Theorem~\ref{thm:rs}), it counts the permutations $\sigma\in\mathfrak S_a$ whose tableau $P(\sigma)$ has a first row of length $\le n$, and by Schensted's theorem (Theorem~\ref{thm:schensted}) this length is the longest
increasing subsequence. If $a\le n$, every permutation qualifies.
\end{proof}

Ira Gessel had found in 1990 the generating function of permutations by longest increasing subsequence as a Toeplitz determinant, which is what this theorem says once Weyl's formula is written out; Baik, Deift and Johansson started from that identity in 1999
to find the Tracy--Widom law of the beauty box of Chapter~9. Ulam's problem about random permutations and the moments of random unitary matrices are the same problem.

\begin{corollary}\label{cor:gaussiantrace}
As $n\to\infty$, $\sqrt2\operatorname{Re}\operatorname{tr}U$ converges in law to the standard Gaussian $N(0,1)$, and so does $\sqrt2\operatorname{Re}(e^{-i\phi}\operatorname{tr}U)$ for every fixed $\phi$.
\end{corollary}

\begin{proof}
The law of $\operatorname{tr}U$ is invariant under rotation (as in the proof of Proposition~\ref{prop:powertraces}), so it suffices to take $\phi=0$; the same argument gives $\mathbb E[(\operatorname{tr}U)^j(\overline{\operatorname{tr}U})^l]=0$ for $j\ne l$. Write $X=\operatorname{tr}U$ and
$\operatorname{Re}X=\frac12(X+\bar X)$. For $n\ge m$, expanding with the binomial formula and Theorem~\ref{thm:traces},
\[
\mathbb E\bigl[(\operatorname{Re}X)^m\bigr]=\frac1{2^m}\sum_{j=0}^m\binom mj\mathbb E\bigl[X^j\bar X^{m-j}\bigr]=\begin{cases}\frac1{2^m}\binom m{m/2}\bigl(\frac m2\bigr)!&m\text{ even},\\0&m\text{ odd}.\end{cases}
\]
For $m=2r$ this is $\frac{(2r)!}{4^rr!}$, so $\mathbb E[(\sqrt2\operatorname{Re}X)^{2r}]=\frac{(2r)!}{2^rr!}$, the moments of $N(0,1)$. The moments of $\sqrt2\operatorname{Re}\operatorname{tr}U$ are thus eventually equal to those of $N(0,1)$, and Tome~5, Lemma~9.6, concludes.
\end{proof}

The trace of a random unitary matrix of size $1000$ is not near $0$ with an error of size $\sqrt{1000}$, as a sum of $1000$ independent unit complex numbers would be: it is a standard complex Gaussian, of size $1$, whatever the size of the matrix.

\begin{beauty}
In 2000 Jon Keating and Nina Snaith computed, with Weyl's integration formula, the moments $\mathbb E|\det(I-U)|^{2k}$ of the characteristic polynomial of a random unitary matrix at the point $1$, and proposed that, in the right scaling, they predict the moments of the
Riemann zeta function on the critical line, $\frac1T\int_0^T|\zeta(\frac12+it)|^{2k}\dd t$. Their formula reproduces the theorems of Hardy and Littlewood ($k=1$, 1918) and Ingham ($k=2$, 1926), and gives for $k=3$ and $k=4$ the constants $42$ and $24024$ that number theorists
(Conrey and Ghosh, Conrey and Gonek) had conjectured by entirely different heuristics. No one knows why it works; it is the strongest evidence that the zeros of $\zeta$ behave like the eigenvalues of a random unitary matrix, the
story that Tome~8, Chapter~8, began with Montgomery and Dyson. Exercise~\ref{ex:charpoly} computes the case $k=1$.
\end{beauty}

\section{Exercises}

\begin{exercise}[$\star$]\label{ex:dims}
Use Theorem~\ref{thm:weylchar} to compute the dimensions of $\operatorname{Sym}^k\C^n$, $\Lambda^k\C^n$ and $V_{(1,0,\dots,0,-1)}$, and identify the last one inside $\C^n\otimes\overline{\C^n}\cong\operatorname{End}(\C^n)$.
\end{exercise}

\begin{exercise}[$\star$]\label{ex:tensorsquare}
Let $n\ge2$.
Show that $s_{(1,0,\dots,0)}^2=s_{(2,0,\dots,0)}+s_{(1,1,0,\dots,0)}$ and deduce that $\C^n\otimes\C^n\cong\operatorname{Sym}^2\C^n\oplus\Lambda^2\C^n$ as representations of $U(n)$, both summands being irreducible.
\end{exercise}

\begin{exercise}[$\star\star$]\label{ex:charpoly}
Show that $\det(I-U)=\sum_{k=0}^n(-1)^k\chi_{\Lambda^k}(U)$ and deduce that $\mathbb E|\det(I-U)|^2=n+1$ for a Haar random $U\in U(n)$.
\end{exercise}

\begin{exercise}[$\star\star$]\label{ex:su2}
(a) Show that the restriction of $V_{(m,0)}$ to $SU(2)$ has the character $\chi_m=\frac{\sin(m+1)\alpha}{\sin\alpha}$ at the matrix with eigenvalues $e^{\pm i\alpha}$. (b) Using the density $\frac2\pi\sin^2\alpha$ of Tome~8, Exercise~8.3, show that these restrictions are irreducible and
pairwise non-isomorphic, and that every irreducible representation of $SU(2)$ is one of them.
\end{exercise}

\begin{exercise}[$\star\star$]\label{ex:catalanu2}
For $U$ Haar distributed in $U(2)$, compute $\mathbb E|\operatorname{tr}U|^{2a}$ directly with Weyl's integration formula, and deduce from Theorem~\ref{thm:traces} that the number of permutations of $\{1,\dots,a\}$ with no increasing subsequence of length $3$ is the
Catalan number $\frac1{a+1}\binom{2a}a$.
\end{exercise}

\begin{exercise}[$\star\star\star$]\label{ex:domino}
Show that $p_2s_\lambda=\sum_\nu\pm s_\nu$, where $\nu$ runs over the dominant weights obtained from $\lambda$ by adding a \emph{domino}: two boxes in the same row (sign $+$) or in the same column (sign $-$) of the diagram, whose union with $\lambda$ is a diagram. (This is the
case $k=2$ of the Murnaghan--Nakayama rule, used in the notebook of Chapter~10.)
\end{exercise}

\section{Hints}

\noindent\textbf{\ref{ex:dims}.} $\delta_j-\delta_k=k-j$. For the last, the trace is an invariant linear form on $\operatorname{End}(\C^n)$.

\noindent\textbf{\ref{ex:tensorsquare}.} Lemma~\ref{lem:pieri}; then compare characters.

\noindent\textbf{\ref{ex:charpoly}.} Expand $\prod_j(1-t_j)$.

\noindent\textbf{\ref{ex:su2}.} For (a), $s_{(m,0)}(t_1,t_2)=\frac{t_1^{m+1}-t_2^{m+1}}{t_1-t_2}$. For completeness, a class function on $SU(2)$ is a function of $\alpha$, and $g(\alpha)\sin\alpha$ extends to an odd function.

\noindent\textbf{\ref{ex:catalanu2}.} $|t_1+t_2|^2|t_1-t_2|^2$ depends only on $\varphi=\theta_1-\theta_2$; use $\int_0^\pi\cos^{2a}x\dd x=\pi\binom{2a}a4^{-a}$.

\noindent\textbf{\ref{ex:domino}.} As in Lemma~\ref{lem:pieri}, $p_2a_{\lambda+\delta}=\sum_ia_{\lambda+\delta+2e_i}$; sort the entries of $\lambda+\delta+2e_i$.

\section{Solutions}

\begin{solution}{ex:dims}
\idea Each factor of the product formula compares two gaps.
\writeup For $\lambda=(k,0,\dots,0)$, only pairs with $j=1$ contribute, and $\dim=\prod_{k'=2}^n\frac{k+k'-1}{k'-1}=\frac{(k+1)(k+2)\cdots(k+n-1)}{(n-1)!}=\binom{n+k-1}k$, the number of monomials of degree $k$ in $n$ variables. For $\lambda=(1^k,0^{n-k})$, the factor for
$j\le k<l$ is $\frac{l-j+1}{l-j}$ and the others are $1$; for fixed $j$, the product over $l=k+1,\dots,n$ telescopes to $\frac{n-j+1}{k-j+1}$, and the product over $j=1,\dots,k$ is $\frac{n(n-1)\cdots(n-k+1)}{k!}=\binom nk$. For $\lambda=(1,0,\dots,0,-1)$, the
factors are $\frac{l}{l-1}$ for $(1,l)$ with $1<l<n$, $\frac{n-j+1}{n-j}$ for $(j,n)$ with $1<j<n$, and $\frac{n+1}{n-1}$ for $(1,n)$; the two telescoping products are each $n-1$, so $\dim=(n-1)^2\frac{n+1}{n-1}=n^2-1$. The character of $\operatorname{End}(\C^n)=\C^n\otimes\overline{\C^n}$ is
$|t_1+\dots+t_n|^2=n+\sum_{j\ne k}t_jt_k^{-1}$, and the trivial representation (the scalars $\C I$, the multiples of the invariant form $\operatorname{tr}$) splits off; what remains is the space of trace-zero matrices, of dimension $n^2-1$, whose character
$n-1+\sum_{j\ne k}t_jt_k^{-1}$ has largest monomial $t_1t_n^{-1}$. Decomposing it into irreducibles $V_\nu$, the largest monomial of the sum is the largest of the leading monomials $t^\nu$ of the constituents (Lemma~\ref{lem:alternants}(a); coefficients are non-negative, so nothing cancels), so $V_{(1,0,\dots,0,-1)}$ is a constituent, and having dimension $n^2-1$ it is everything: the trace-zero matrices form $V_{(1,0,\dots,0,-1)}$, the \emph{adjoint representation} on $\mathfrak{sl}_n(\C)$.
\end{solution}

\begin{solution}{ex:tensorsquare}
\idea One box added to one box, in two ways; then count irreducible pieces.
\writeup By Lemma~\ref{lem:pieri} with $\lambda=(1,0,\dots,0)$, the dominant weights obtained by adding $1$ to one entry are $(2,0,\dots,0)$ and $(1,1,0,\dots,0)$, so $p_1^2=s_{(2)}+s_{(1,1)}$ (for $n\ge2$). The character of $\C^n\otimes\C^n$ is $p_1^2$, so by Theorem~\ref{thm:schurorth}
$\C^n\otimes\C^n\cong V_{(2)}\oplus V_{(1,1)}$, a sum of exactly two irreducibles. The symmetric and antisymmetric tensors are non-zero invariant subspaces with $\C^n\otimes\C^n=\operatorname{Sym}^2\oplus\Lambda^2$; each is a sum of irreducibles, and there are only two in total, so each is
irreducible. The weight $(2,0,\dots,0)$, that is the monomial $t_1^2$, occurs in $\operatorname{Sym}^2$ (on $e_1\otimes e_1$) and not in $\Lambda^2$, and it occurs in $s_{(2)}$ but not in $s_{(1,1)}=e_2$; so $\operatorname{Sym}^2\C^n\cong V_{(2)}$ and $\Lambda^2\C^n\cong V_{(1,1)}$.
\end{solution}

\begin{solution}{ex:charpoly}
\idea The characteristic polynomial at $1$ is an alternating sum of irreducible characters.
\writeup On $T$, $\det(I-U)=\prod_j(1-t_j)=\sum_k(-1)^ke_k(t)$, and $e_k=s_{(1^k,0^{n-k})}$ is the character of $\Lambda^k\C^n$, irreducible by Theorem~\ref{thm:weylchar}. By orthonormality of the $n+1$ distinct characters $e_0,\dots,e_n$,
$\mathbb E|\det(I-U)|^2=\sum_{k=0}^n|(-1)^k|^2=n+1$. (For $k=1$ the Keating--Snaith formula is $\prod_{j=0}^{n-1}\frac{\Gamma(j+1)\Gamma(j+3)}{\Gamma(j+2)^2}=\prod_{j=0}^{n-1}\frac{j+2}{j+1}=n+1$.)
\end{solution}

\begin{solution}{ex:su2}
\idea Restrict Weyl's formula from $U(2)$ to $SU(2)$, and use the Fourier sine series on $[0,\pi]$ for completeness.
\writeup (a) By Theorem~\ref{thm:weylchar}, with $\delta=(1,0)$, $s_{(m,0)}=\frac{t_1^{m+1}t_2^0-t_2^{m+1}t_1^0}{t_1-t_2}$; at $t_1=e^{i\alpha}$, $t_2=e^{-i\alpha}$ this is $\frac{2i\sin(m+1)\alpha}{2i\sin\alpha}$.
(b) A class function on $SU(2)$ depends only on $\alpha\in[0,\pi]$ (eigenvalues $e^{\pm i\alpha}$ up to order), and its integral is $\int_0^\pi g(\alpha)\frac2\pi\sin^2\alpha\dd\alpha$. Then $\langle\chi_m,\chi_{m'}\rangle=\frac2\pi\int_0^\pi\sin(m+1)\alpha\sin(m'+1)\alpha\dd\alpha=\delta_{mm'}$:
so the restrictions are irreducible (norm $1$, Theorem~\ref{thm:schurorth}) and pairwise non-isomorphic. If an irreducible $\rho$ of $SU(2)$ were none of them, its character $g(\alpha)$ would satisfy $\int_0^\pi g(\alpha)\sin\alpha\,\sin(m+1)\alpha\dd\alpha=0$ for all $m\ge0$. The function
$g(\alpha)\sin\alpha$ is continuous on $[0,\pi]$, and its odd extension to $[-\pi,\pi]$ has all its Fourier coefficients zero (the cosine coefficients vanish by oddness), so it is $0$ (Theorem~\ref{thm:peterweyl}(b) for $U(1)$: a continuous function orthogonal to every $e^{ik\theta}$ vanishes): $g=0$ on $]0,\pi[$, hence everywhere by continuity, contradicting $\langle\chi_\rho,\chi_\rho\rangle=1$.
\end{solution}

\begin{solution}{ex:catalanu2}
\idea Weyl's formula for $n=2$ reduces the moment to a one-variable integral, and Theorem~\ref{thm:traces} reads it as a count.
\writeup With $t_j=e^{i\theta_j}$ and $\varphi=\theta_1-\theta_2$: $|t_1+t_2|^2=2+2\cos\varphi=4\cos^2\frac\varphi2$ and $|t_1-t_2|^2=2-2\cos\varphi=4\sin^2\frac\varphi2$. Weyl's formula gives
\[
\mathbb E|\operatorname{tr}U|^{2a}=\frac12\cdot\frac1{2\pi}\int_0^{2\pi}4^a\cos^{2a}\frac\varphi2\cdot4\sin^2\frac\varphi2\dd\varphi=\frac{2\cdot4^a}\pi\int_0^\pi\cos^{2a}x\,(1-\cos^2x)\dd x
\]
(the double integral over $\theta_1,\theta_2$ reduces to one over $\varphi$, by periodicity). With $I_a=\int_0^\pi\cos^{2a}x\dd x=\pi\binom{2a}a4^{-a}$ (Wallis), and $\binom{2a+2}{a+1}=\binom{2a}a\frac{2(2a+1)}{a+1}$, this is
$2\binom{2a}a\bigl(1-\frac{2a+1}{2a+2}\bigr)=\frac1{a+1}\binom{2a}a$. By Theorem~\ref{thm:traces} with $n=2$, it is also the number of $\sigma\in\mathfrak S_a$ whose longest increasing subsequence has length $\le2$.
\end{solution}

\begin{solution}{ex:domino}
\idea Adding $2$ to one entry of the strictly decreasing vector $l=\lambda+\delta$ either collides with the next entry, or jumps over it, or neither; each jump is a sign.
\writeup As in the proof of Proposition~\ref{prop:powertraces}, $p_2a_l=\sum_ia_{l+2e_i}$, with $l=\lambda+\delta$ strictly decreasing. Fix $i$ and look at the entry $l_i+2$. Since the entries are distinct integers, $l_{i-1}\ge l_i+1$. If $l_{i-1}=l_i+2$, two entries coincide and the
term vanishes. If $l_{i-1}\ge l_i+3$ (or $i=1$), the vector $l+2e_i$ is strictly decreasing, equal to $\nu+\delta$ with $\nu=\lambda+2e_i$: two boxes added to row $i$, a horizontal domino, sign $+$. If $l_{i-1}=l_i+1$, then $l_i+2=l_{i-1}+1$; exchanging positions $i-1$ and $i$ (a
sign $-1$) gives the vector with entries $l_i+2$ in position $i-1$ and $l_{i-1}$ in position $i$, which is strictly decreasing if $l_{i-2}\ne l_i+2$, that is $l_{i-2}\ne l_{i-1}+1$ (if equal, the term vanishes; it cannot be smaller, as $l_{i-2}>l_{i-1}$). In terms of $\lambda$: $l_{i-1}=l_i+1$
means $\lambda_{i-1}=\lambda_i$, and the new weight is $\nu$ with $\nu_{i-1}=\lambda_{i-1}+1$, $\nu_i=\lambda_i+1$: one box added at the end of row $i-1$ and one below it at the end of row $i$, a vertical domino, sign $-$; it is a diagram exactly when $\lambda_{i-2}>\lambda_{i-1}$,
the condition $l_{i-2}\ne l_{i-1}+1$. Every domino addition arises once in this way. Dividing by $a_\delta$ gives the rule.
\end{solution}

\section*{Sources and further reading}
\addcontentsline{toc}{section}{Sources and further reading}

I.~Schur, ``Neue Anwendungen der Integralrechnung auf Probleme der Invariantentheorie'', \emph{Sitzungsber.\ Preuss.\ Akad.\ Wiss.} (1924), 189--208, 297--321, 346--355; H.~Weyl, ``Theorie der Darstellung kontinuierlicher halb-einfacher Gruppen durch lineare
Transformationen'', \emph{Math.\ Z.} 23 (1925), 271--309, and 24 (1926), 328--395, 789--791; F.~Peter and H.~Weyl, ``Die Vollständigkeit der primitiven Darstellungen einer geschlossenen kontinuierlichen Gruppe'', \emph{Math.\ Ann.} 97 (1927), 737--755;
H.~Weyl, ``Invariants'', \emph{Duke Math.\ J.} 5 (1939), 489--502, for the quotation; H.~Weyl, \emph{The Classical Groups} (Princeton University Press, 1939). P.~Diaconis and M.~Shahshahani, ``On the eigenvalues of random matrices'', \emph{J.\ Appl.\ Probab.} 31A (1994), 49--62;
E.~M.~Rains, ``Increasing subsequences and the classical groups'', \emph{Electron.\ J.\ Combin.} 5 (1998), R12; I.~M.~Gessel, ``Symmetric functions and P-recursiveness'', \emph{J.\ Combin.\ Theory Ser.\ A} 53 (1990), 257--285; J.~P.~Keating and N.~C.~Snaith, ``Random matrix
theory and $\zeta(1/2+it)$'', \emph{Comm.\ Math.\ Phys.} 214 (2000), 57--89; M.~Gell-Mann, \emph{The Eightfold Way: a Theory of Strong Interaction Symmetry}, Caltech report CTSL-20 (1961).

For more: T.~Bröcker and T.~tom Dieck, \emph{Representations of Compact Lie Groups} (Springer, 1985), Chapters~II--VI; D.~Bump, \emph{Lie Groups} (Springer, 2nd edition 2013), for Schur--Weyl duality and the Keating--Snaith formula; I.~G.~Macdonald, \emph{Symmetric Functions and
Hall Polynomials} (Oxford, 2nd edition 1995), Chapter~I; E.~S.~Meckes, \emph{The Random Matrix Theory of the Classical Compact Groups} (Cambridge University Press, 2019).
